%% kapa: exam=1 round=33 year=2022 period=1 form=A subject=civil_law subject_ko=민법 seq=19 number=19
%% source: Q-Net 감정평가사 기출문제 (cst00309) artlSeq=5209997 file=raw/1008/5209997/2022년도 제33회 감정평가사 1차시험 1교시(민법,경제학,부동산학).pdf
%% conversion: kordoc 4.24.51 -> tools/gen1.py (auto); review=proofread-sonnet; flags=none
\begin{question}{19}
소멸시효의 기산점이 옳게 연결되지 않은 것은? (다툼이 있으면 판례에 따름)
\begin{choices}
  \item 부작위를 목적으로 하는 채권 - 위반행위시
  \item 동시이행의 항변권이 붙어 있는 채권 - 이행기 도래시
  \item 이행불능으로 인한 손해배상청구권 - 이행불능시
  \item 甲이 자기 소유의 건물 매도시 그 이익을 乙과 분배하기로 약정한 경우 乙의 이익금 분배청구권 - 분배약정시
  \item 기한이 있는 채권의 이행기가 도래한 후 채권자와 채무자가 기한을 유예하기로 합의한 경우 그 채권 - 변경된 이행기 도래시
\end{choices}
\end{question}
\answer{4}
